\section{Discussion}
\label{sec:discussion}

\subsection{Labels and validation decide the accuracy figure}

The label source is the result that bears most directly on published flood maps. A model trained on terrain-threshold labels and scored against observed flooding on the same held-out blocks lost \RRLabelGap{} in AUC in Rangpur and Rajshahi, \SylhetLabelGap{} in Sylhet and \CoastalLabelGap{} on the coast, trailing on \RRLabelAhead{}, \SylhetLabelAhead{} and \CoastalLabelAhead{} of the \RRBenchSeeds{} block assignments respectively. The wetness index on its own, a common predictor in the studies of Table~\ref{tab:litflood}, ranked flood-prone assets at \RRBenchTwiAUC{} in Rangpur and Rajshahi and \SylhetBenchTwiAUC{} in Sylhet, close to chance, which means the signal lies in the combination of features rather than in any single terrain threshold. A model scored against labels derived from the same terrain it takes as input can succeed by relearning the labelling rule, and its score then measures recovery of that rule rather than skill at locating floods.

The AUCs reported here, between \RRObsAUC{} and \SylhetObsAUC{} for the riverine regions, sit below the 0.88 to 0.96 reported by the studies in Table~\ref{tab:litflood}. The comparison is only indicative, since those studies differ in extent, inventory and model, and their splits are mostly not stated. The gap need not mean a weaker model, because every figure here comes from 10\,km blocks the model never saw, scored against flooding the radar recorded, and is averaged over several block assignments whose spread, between \RRBenchBoostingSD{} and \CoastalBenchBoostingSD{} in AUC, is as large as several of the differences reported here. Holding out whole areas has its own critics, who argue that it can force a model to extrapolate and so overstate error \citep{roberts2017}, and that it is no substitute for a probability sample of the finished map \citep{wadoux2021}. The blocked scores therefore estimate performance on unseen areas of similar terrain, and they are not design-based measures of map accuracy. \pending{extend this paragraph once the literature has been checked against full texts.}

\subsection{What the graph did not add}

A graph neural network was the starting design of this system, on the reasoning that neighbouring assets flood together and a model that shares information between them should rank them better. On identical features, assets and blocks it trailed the best ordinary model by \RRBenchGap{} in AUC in Rangpur and Rajshahi ($p\RRBenchDiffP{}$) and \SylhetBenchGap{} in Sylhet ($p\SylhetBenchDiffP{}$), and by \CoastalBenchGap{} on the coast, where the difference is within noise ($p\CoastalBenchDiffP{}$). A likely reason, not tested here, is that the features already carry the spatial information the graph was meant to supply, since assets that are close together share terrain, wetness and distances to the same roads and shelters, and a tree can split on those values directly. Interpolating the model's scores by Kriging failed for a related reason, with the terrain surface scored cell by cell outperforming the Kriged surface against observed flooding in all three regions (\RRSurfaceTerrainObsAUC{} against \RRSurfaceKrigedObsAUC{} in Rangpur and Rajshahi), because Kriging smooths the scores over tens of kilometres while flooding follows individual channels and levees. At this scale and with these features, neither linking assets into a graph nor interpolating between them improved spatial prediction, and studies proposing either should report the tabular baseline on the same blocks.

\subsection{What the flood record already knows}

The temporal hold-out adds a test that susceptibility studies rarely report, which is whether a model anticipates a later flood better than the record of earlier ones. On its own the model did not, in either riverine region, with the share of earlier events in which the ground flooded ranking the 2024 flooding at \RRTempPastAUC{} and \SylhetTempPastAUC{} against \RRTempAUC{} and \SylhetTempAUC{} for the model (Section~\ref{sec:temporal}). Given that record as a feature, the model beat both, at \RRPastFeatAUC{} and \SylhetPastFeatAUC{}, and it did so on the ground no earlier flood had reached, where the record alone cannot rank at all. The two sources therefore carry different information, with the flood record marking where water has returned and the terrain ranking where it has not yet been seen, and a susceptibility map for a river basin with a radar flood history should be built on both. On the coast, where \CoastalPastFeatUnseenPos{}\,\% of the assets Cyclone Remal flooded stood on ground no earlier cyclone had reached, the model trained on all three earlier cyclones did best, at \CoastalTempAUC{} against \CoastalTempPastAUC{} for the record alone. A study that reports accuracy without the past-flooding reference cannot show that its model knows more than the flood record already does, and the two are simple to compare wherever Sentinel-1 has observed earlier floods.

\subsection{Where the maps should not be trusted}

The coast is the clearest limit. There the asset model still ranks flooded ground above dry ground, but neither hazard surface does better than chance across the grid, because surge flooding is governed by polder embankments, tidal channels and landfall position, none of which the elevation-derived predictors describe. A coastal model would need embankment and tide data that the pipeline does not yet ingest, and until then the coastal composite should not be read as a map of surge exposure.

Exposure depends on OpenStreetMap, whose completeness across rural Bangladesh has not been measured, and composite risk is defined only where assets are mapped, which in Rangpur and Rajshahi is 7\,\% of grid cells. Units without mapped assets are therefore left undefined, since sparse mapping would otherwise read as low risk. The infrastructure composite rises with asset density and is led by urban municipalities, so it measures infrastructure at risk, and risk to people is carried by the population composite. Vulnerability rests on population density and three access distances, so income, housing quality and age structure are not represented, and WorldPop's constrained product misses settlement outside detected buildings.

Two caveats concern the flood labels themselves. Inundated paddy is open water to radar and is counted as flooding, which is one reason the riverine labels require flooding in at least two events, and the minimum-backscatter composite over a two-week window records the largest extent reached at any pass in the window. The calibrated probabilities approximate likelihoods for a region as a whole and understate them where flood-prone ground is concentrated, as the shift in prevalence between calibration and test blocks showed (Fig.~\ref{fig:calibration}). None of the surfaces is a forecast, the scores describe typical monsoon conditions and give no probability of flooding in a given year, and warnings remain the task of the Flood Forecasting and Warning Centre and the Bangladesh Meteorological Department. The landslide surface describes which slopes failed in one storm, and retraining on the multi-year inventory of \citet{rabby2019} is the most direct improvement available.

\section{Conclusions}
\label{sec:conclusions}

Training on flood extents observed by Sentinel-1 raised the AUC over terrain-threshold labels by \RRLabelGap{} in Rangpur and Rajshahi, \SylhetLabelGap{} in Sylhet and \CoastalLabelGap{} on the coast, on ground the models had never seen. This is the finding with the widest consequence, because it bears on how any susceptibility map trained and scored against terrain-derived labels should be read. Ordinary tabular models beat a graph neural network in every region, measurably so in the two riverine ones, so the features rather than the architecture carried the signal. The terrain hazard surface followed observed flooding in the riverine regions and failed on the coast, where surge depends on structures that elevation does not record. Scored against the 2024 floods after training on earlier ones, the model held most of its skill, and in the riverine regions a model given the record of earlier flooding as a feature outperformed both that record and the model without it, reaching \RRPastFeatAUC{} and \SylhetPastFeatAUC{}. Every figure in this paper is generated from the pipeline outputs, and the data, code and dashboard are released so that each can be checked and rebuilt.

